Introduction goes here

% TODO: Introduction should describe the current solutions for the problem, and
% why there's a gap between the needs and solutions that needs to be covered.

This is an example of how to cite a reference
\parencite{digitalization-europe-goals}. You also have the option of citing the
author \citeauthor{digitalization-europe-goals}, or the year
\citeyear{digitalization-europe-goals}. Or both if you use textcite
\textcite{digitalization-europe-goals}. References are stored in references.bib

You also have the option to use acronyms, which are stored in acronyms.tex. You
can use the short version \acrshort{hci} or the longer one \acrlong{hci} or
full version \acrfull{hci}. You even have the option to configure plurals,
depending on whether they are different from singulars: singular short
\acrshort{sme}, plural short \acrshortpl{sme}, singular long \acrlong{sme},
singular full \acrfull{sme}, plural full \acrfullpl{sme}. These are all
clickable, just like the references, and will take you to the full details in
the glossary if you click on them.

You can also add images, by the way. To follow all the standards from the uni,
I created a custom command called addfigure. You can also reference them if
needed by using the command ref: Figure \ref{fig:sample_figure}.

\addfigure{Example title on how to add a figure}{assets/Picture1.jpg}{caption for the
	image}{fig:sample_figure}

Of course, you also have the option to add tables. In Table
\ref{tbl:table_example} you can find an example of how to create one.

\newpage

\addtable{This is an example on how to create a table}{%
	\begin{tabular}{|l|l|}
		\hline
		One   & Two  \\
		\hline
		Three & Four \\
		\hline
	\end{tabular}
}{Caption for the table}{tbl:table_example}

Both tables and figures will be added to the corresponding directories, which
can be found after the index.
